\documentclass{article}
% ============================================================
% Bibliography and citations
% ============================================================

\usepackage{natbib}
\usepackage{csquotes}


% ============================================================
% Mathematics
% ============================================================

\usepackage{
    amsmath,
    amssymb,
    amsthm,
    amsfonts,
    mathrsfs,
    dsfont,
    bm,
    textgreek,
    bbm
}


% ============================================================
% Fonts
% ============================================================

% \usepackage{newpxtext}
% \usepackage{eulervm}


% ============================================================
% Page layout and document structure
% ============================================================

\usepackage[margin=1in, footskip=0.3in]{geometry}
\usepackage{setspace}
\usepackage[bf,small]{titlesec}
\usepackage{etoc}
\usepackage{chngcntr}
\usepackage{changepage}
\usepackage{afterpage}


% ============================================================
% Figures, captions, and graphics
% ============================================================

\usepackage{graphicx}
\usepackage[bf,skip=8pt]{caption}
\usepackage{subcaption}
\usepackage{tikz}

\graphicspath{{Figures/}}


% ============================================================
% Tables
% ============================================================

\usepackage{booktabs}
\usepackage{tabularx}
\usepackage{dcolumn}
\usepackage{makecell}
\usepackage{multirow}
\usepackage{longtable}
\usepackage[para]{threeparttable}
\usepackage{colortbl}
\usepackage{adjustbox}


% ============================================================
% Landscape and rotated content
% ============================================================

\usepackage{lscape}
\usepackage{pdflscape}
\usepackage{rotating}


% ============================================================
% Lists, comments, notes, and text utilities
% ============================================================

\usepackage{enumerate}
\usepackage{comment}
\usepackage{verbatim}
\usepackage{soul}
\usepackage[multiple,stable]{footmisc}
\usepackage[colorinlistoftodos]{todonotes}


% ============================================================
% Code listings
% ============================================================

\usepackage{listings}


% ============================================================
% Colors
% ============================================================

\usepackage{xcolor}

% General-purpose colors
\definecolor{blue}{rgb}{0,0.08,0.5}
\definecolor{red}{rgb}{0.6,0,0}
\definecolor{green}{rgb}{0,0.5,0.05}

% Code-listing colors
\definecolor{fgreen}{HTML}{698B27}
\definecolor{lavender}{HTML}{E6E6FA}
\definecolor{darkpink}{HTML}{C71585}
\definecolor{darkyellow}{HTML}{B8860B}


% ============================================================
% Hyperlinks and PDF metadata
% ============================================================

\usepackage[
    bookmarks,
    pdftitle={},
    pdfauthor={Banks},
    colorlinks=true,
    linkcolor=black,
    citecolor=black,
    urlcolor=black
]{hyperref}


% ============================================================
% Theorem environments
% ============================================================

\newtheorem{prediction}{Prediction}
\newtheorem{assumption}{Assumption}
\newtheorem{proposition}{Proposition}
\newtheorem{corollary}{Corollary}
\newtheorem{definition}{Definition}
\newtheorem{lemma}{Lemma}

% Proposition numbering resets by section
\newtheorem{prop}{Proposition}[section]


% ============================================================
% Mathematical commands
% ============================================================

% Formatting
\newcommand{\no}{\noindent}
\newcommand{\mb}{\mathbf}

% Expectation, probability, and covariance
\newcommand{\E}{\mathbb{E}}
\newcommand{\R}{\mathbb{R}}
\newcommand{\pr}{\mathrm{Pr}}
\newcommand{\cov}{\mathrm{Cov}}

% Support
\newcommand{\supp}[1]{\mathrm{Supp}\left(#1\right)}

% Partial derivatives
\newcommand{\pd}[2]{\frac{\partial #1}{\partial #2}}
\newcommand{\pdd}[2]{\frac{\partial^2 #1}{\partial #2^2}}

% Indicator function
\newcommand{\ind}[1]{\mathbbm{1}\{#1\}}

% Vertical spacing helper for large fractions
\newcommand{\vh}{
    \vphantom{
        \frac{
            \frac{\pd{a}{a}}{\pd{a}{a}}
        }{
            \frac{\pd{a}{a}}{\pd{a}{a}}
        }
    }
}


% ============================================================
% Custom environments
% ============================================================

% Wide equation extending 2 cm into each margin
\newenvironment{widerequation}
{
    \begin{adjustwidth}{-2cm}{-2cm}
    \begin{equation*}
}
{
    \end{equation*}
    \end{adjustwidth}
}

% Solution environment
\newenvironment{solution}
{
    \par\noindent\textbf{\textit{Solution.}}
}
{
    \hfill\rule{1.1ex}{1.1ex}
}


% ============================================================
% Code-listing style
% ============================================================

\lstset{
    basicstyle=\ttfamily\small\color{fgreen},
    keywordstyle=\color{darkpink},
    commentstyle=\color{gray},
    stringstyle=\color{darkyellow},
    showstringspaces=false,
    breaklines=true,
    frame=single,
    backgroundcolor=\color{lavender}
}


% ============================================================
% Spacing and section formatting
% ============================================================

% ECMA-style line spacing
\linespread{1.5}

\titlespacing*{\paragraph}
    {0pt}
    {.25\baselineskip}
    {1em}

\titlespacing*{\section}
    {0pt}
    {.6\baselineskip}
    {.3\baselineskip}

\titlespacing*{\subsection}
    {0pt}
    {.6\baselineskip}
    {.3\baselineskip}

\titlespacing*{\subsubsection}
    {0pt}
    {.6\baselineskip}
    {.3\baselineskip}

% \setlength{\parskip}{5pt}

\setlength{\footnotesep}{.2in}


% ============================================================
% Subsubsection formatting
% ============================================================

\titleformat{\subsubsection}[runin]
    {\normalfont\itshape\bfseries}
    {\thesubsubsection}
    {1em}
    {}
    [.]


% ============================================================
% Date and title-page utilities
% ============================================================

\newcommand{\monthyear}{%
    \ifcase\month\or
        January\or
        February\or
        March\or
        April\or
        May\or
        June\or
        July\or
        August\or
        September\or
        October\or
        November\or
        December
    \fi
    \ \number\year
}


\begin{document}
\begin{center}
    \textbf{\scalebox{1.3}{Economics 880 PSet 1  (Lentz)}}\\
    Aaron Banks \& Arthur Rodrigues
\end{center}

\section*{Question 1}

Following the defintion of $\ln L(Y_i \mid \theta)$ given in setting of the problem set, we have 
\begin{equation}\label{eq:q1_L}
    \begin{aligned}
        \mathcal{L}(Y \mid \hat{\theta}) &= \sum_{i=1}^I \ln L(Y_i \mid \hat{\theta}) \\ 
        &= \sum_{i=1}^I \sum_{s=1}^{T_i}\left(
            - \frac{(\ln \Delta x_{is} - \hat{\mu})^2}{2\hat{\sigma}^2} - 
            \ln(\Delta x_{is}) - \ln \hat{\sigma} - \frac{1}{2}\ln(2\pi)
        \right)\\
        &\quad + \quad
        \sum_{i=1}^I \sum_{s=1}^{T_i}\sum_{n=1}^N 
        \chi_{isn} [
            d_{ist+1}\ln \hat{h}_n + (1-d_{ist+1})\ln (1 - \hat{h}_n)
        ].
    \end{aligned}
\end{equation}
The value $\hat{\theta}$ maximizes $\mathcal{L}(Y \mid \hat{\theta})$, so we can use 
the first order condition of the latter to derive the components of $\hat{\theta}$. 
For $\hat{\mu}$, we have 
\begin{align*}
    \frac{\partial \mathcal{L}}{\partial \hat{\mu}}
     = \sum_{i=1}^I \sum_{s=1}^{T_i}\left(
            - \frac{2 (\ln \Delta x_{is} - \hat{\mu})}{2 \hat{\sigma}^2} \cdot (-1)
        \right)
        \quad \overset{FOC}{\Rightarrow} \quad 
        \hat{\mu} = 
        \frac{1}{\sum_{i=1}^I T_i} \sum_{i=1}^I \sum_{s=1}^{T_i} \ln \Delta x_{is}.
\end{align*}
For $\hat{\sigma}$, we have
\begin{align*}
    \frac{\partial \mathcal{L}}{\partial \hat{\sigma}} = 
    \sum_{i=1}^I\sum_{s}^{T_i} \left(
        (\ln \Delta x_{is} - \hat{\mu})^2 (2 \hat{\sigma}^2)^{-2}4\sigma - \sigma^{-1}
    \right)
    \quad \overset{FOC}{\Rightarrow} \quad 
    \hat{\sigma}^2 = \frac{1}{ \sum_{i=1}^I T_i} \sum_{i=1}^I\sum_{s=1}^{T_i} \left(
        \ln \Delta x_{is} - \hat{\mu}
    \right)^2.
\end{align*}
Finally, for $\hat{h}$, we have 
\begin{align*}
    \frac{\partial \mathcal{L}}{\partial h_j}
    &= \sum_{i=1}^I \sum_{s=1}^{T_i} \chi_{isj}
    \left(
        \frac{d_{is+1}}{h_j}
        - \frac{1-d_{is+1}}{1-h_j}
    \right) 
    \quad \overset{FOC}{\Rightarrow} \quad
    \hat{h}_j
    = \frac{
        \sum_{i=1}^I \sum_{s=1}^{T_i} \chi_{isj}d_{is+1}
    }{
        \sum_{i=1}^I \sum_{s=1}^{T_i} \chi_{isj}
    }.
\end{align*}
Plugging these into $\mathcal{L}$ above yields
\begin{align*}
\mathcal{L}(Y \mid \hat{\theta})
&= -\frac{1}{2}\sum_{i=1}^I T_i
- \sum_{i=1}^I \sum_{s=1}^{T_i} \ln(\Delta x_{is})
- \left(\sum_{i=1}^I T_i\right)
\left(\ln \hat{\sigma} + \frac{1}{2}\ln(2\pi)\right) \\
&\quad + \sum_{j=1}^N \Bigg[
\left(
\sum_{i=1}^I \sum_{s=1}^{T_i}
\chi_{isj}d_{is+1}
\right)
\ln\left(
\frac{
\sum_{i=1}^I \sum_{s=1}^{T_i}
\chi_{isj}d_{is+1}
}{
\sum_{i=1}^I \sum_{s=1}^{T_i}
\chi_{isj}
}
\right) \\
&\qquad\qquad +
\left(
\sum_{i=1}^I \sum_{s=1}^{T_i}
\chi_{isj}(1-d_{is+1})
\right)
\ln\left(
\frac{
\sum_{i=1}^I \sum_{s=1}^{T_i}
\chi_{isj}(1-d_{is+1})
}{
\sum_{i=1}^I \sum_{s=1}^{T_i}
\chi_{isj}
}
\right)
\Bigg].
\end{align*}
However, we will not used the above to compute the likeliehood. Rather, we find it simpler to
first compute $\hat{\theta}$ from the first order conditions above as they define simple 
frequency measures that can be computed directly from the data, and then plugging these into 
equation \eqref{eq:q1_L} to get the likeliehood estimate. We have 
\[
    (\hat{\mu}, \hat{\sigma}, \hat{\mathcal{L}}) = (8.02,0.619,-71146.863).
\]
Figure \ref{fig:q1_hist} plots the CDF of mileage associated with replacement.  

\begin{figure}[ht!]
    \centering
    \caption{Replacement Hazard}
    \label{fig:q1_hist}
    \includegraphics[width=0.75\linewidth]{../output/hazard_cdf.pdf}
\end{figure}

\section*{Question 2}

Rather than estimating the replacement hazard separately over $N=40$ mileage
bins, we now parameterize $h(x)$ using an $n$-order Bernstein polynomial.
In particular, we have
\begin{align*}
    h(x \mid f)
    &= \frac{1}{\bar{x}^n}\sum_{k=1}^n
    f_k \binom{n}{k} x^k(\bar{x}-x)^{n-k} \\
    &= \sum_{k=1}^n f_k B_k(x),
\end{align*}
where
\[
    B_k(x) = \binom{n}{k}
    \left(\frac{x}{\bar{x}}\right)^k
    \left(1-\frac{x}{\bar{x}}\right)^{n-k},
    \qquad f_k \in [0,1].
\]
We set $n=5$ and $\bar{x}=\max_{it}x_{it}$, as specified in the
problem set. The restriction $h(0)=0$ is imposed by omitting the
$k=0$ term.

As in Question 1, the log-likelihood is separable into the odometer
and replacement components. Since the odometer component does not
depend on $f$, we can estimate $(\mu,\sigma)$ analytically as before
and maximize the replacement component with respect to $f$.
We have
\begin{align*}
    R(Y\mid f)
    = \sum_{i=1}^I\sum_{s=1}^{T_i}
    \left[
        d_{i,s+1}\ln h(x_{is}\mid f)
        +(1-d_{i,s+1})\ln(1-h(x_{is}\mid f))
    \right].
\end{align*}
The estimator is therefore given by
\[
    \hat{f} = \operatorname{argmax}_{f\in[0,1]^5} R(Y\mid f).
\]

\begin{enumerate}

\item[(a)]
We first consider estimating $f$ using Newton-Raphson.
For this, we need the gradient and Hessian of the
replacement log-likelihood. Since
\[
    \frac{\partial h(x_{is}\mid f)}{\partial f_j}
    = B_j(x_{is}),
\]
the first derivative is given by
\begin{align*}
    \frac{\partial R}{\partial f_j}
    &= \sum_{i=1}^I\sum_{s=1}^{T_i}
    B_j(x_{is})
    \left[
        \frac{d_{i,s+1}}{h(x_{is}\mid f)}
        -\frac{1-d_{i,s+1}}{1-h(x_{is}\mid f)}
    \right].
\end{align*}
Differentiating again, we have
\begin{align*}
    \frac{\partial^2 R}{\partial f_j\partial f_k}
    &=-\sum_{i=1}^I\sum_{s=1}^{T_i}
    B_j(x_{is})B_k(x_{is})
    \left[
        \frac{d_{i,s+1}}{h(x_{is}\mid f)^2}
        +\frac{1-d_{i,s+1}}{(1-h(x_{is}\mid f))^2}
    \right].
\end{align*}
We use these expressions to implement Newton-Raphson by
minimizing the negative replacement log-likelihood
$Q(f)=-R(Y\mid f)$. The standard Newton-Raphson update
is given by
\[
    f^{(r+1)}
    = f^{(r)}
    - \left[\nabla^2 Q(f^{(r)})\right]^{-1}
    \nabla Q(f^{(r)}).
\]

We first implement the standard Newton-Raphson algorithm
without modifying the step size or imposing the parameter
constraints during the update. Starting from
$f^{(0)}=(0.5,0.5,0.5,0.5,0.5)$, the algorithm fails to
converge because an iteration produces coefficients outside
the feasible parameter space $[0,1]^5$. In particular,
the proposed coefficients at termination are
\[
    f=(-2.6028,\;3.9514,\;-3.6908,\;2.2029,\;-0.6544).
\]
These coefficients are not admissible, and consequently
we do not obtain a valid likelihood estimate.

To examine sensitivity to initial guesses, we repeat the
estimation using
\[
    f^{(0)}\in
    \left\{
    (a,a,a,a,a): a=0.1,0.3,0.5,0.7,0.9
    \right\}.
\]
We find that standard Newton-Raphson fails to converge
for all five initial guesses, with each run eventually
producing a step outside the feasible parameter space.
Although the proposed coefficients at termination differ
across starting values, the reason for failure is the same.
Thus, the standard algorithm fails for all of the initial
guesses considered, rather than converging only from
particular starting values.

We therefore make several modifications to the algorithm
to improve its numerical stability. First, we impose the
parameter restrictions by projecting proposed updates
onto $[0,1]^5$. Second, we use a backtracking line search,
which reduces the step size whenever the proposed update
fails to sufficiently decrease the objective. Third, we
identify coefficients constrained at their boundaries and
exclude them from the Newton system when the gradient points
outside the feasible parameter space. Finally, we regularize
the Hessian by adding a small positive constant to its
diagonal and allow a gradient descent step when the Newton
direction fails to produce an acceptable update. Thus, our
modified algorithm is a damped, constrained
version of Newton-Raphson.

Starting again from $f^{(0)}=(0.5,0.5,0.5,0.5,0.5)$,
the modified algorithm converges to
\[
    \hat{f}_{NR}
    = (0,\;0,\;0.004654,\;0.071000,\;0.037850).
\]
Combining this with the estimates of the odometer parameters,
we obtain
\[
    (\hat{\mu},\hat{\sigma},
    \mathcal{L}(Y\mid\hat{\theta}_{NR}))
    = (8.0165,\;0.6186,\;-71160.713).
\]
The gradient error is $4.289\times10^{-10}$.

Repeating the estimation over the five initial guesses
considered above, we find that the modified algorithm
converges to the same coefficients and maximized
log-likelihood in every case. Thus, the modified
Newton-Raphson algorithm is not particularly sensitive
to the initial guesses considered.

The difference between the two implementations is
particularly important because the estimated first two
Bernstein coefficients are zero, so the maximizer lies
on the boundary of the parameter space. Standard
Newton-Raphson does not account for these restrictions
and can produce infeasible updates, whereas the modified
algorithm accommodates the boundary solution.


    \item[(b)]
    Next, we estimate $f$ using Nelder-Mead. Unlike Newton-Raphson,
    Nelder-Mead does not require the gradient or Hessian of the
    objective. Rather, it maintains a simplex of $n+1$ points and
    updates this simplex using evaluations of the negative
    log-likelihood.

    Using the same initial guess as above, we obtain
    \[
        \hat{f}_{NM}
        = (0,\;0,\;0.004654,\;0.071000,\;0.037850),
    \]
    with
    \[
        \mathcal{L}(Y\mid\hat{\theta}_{NM})
        = -71160.713.
    \]
    Thus, for this initial guess, Nelder-Mead produces the same
    estimates as Newton-Raphson.

    However, Nelder-Mead is considerably more sensitive to initial
    guesses. Repeating the estimation over the same five starting
    values gives the following results:
    \begin{center}
    \begin{tabular}{c c c}
        \hline
        Initial guess & Log-likelihood & Converged \\
        \hline
        0.1 & -71160.774 & No \\
        0.3 & -71160.774 & No \\
        0.5 & -71160.713 & Yes \\
        0.7 & -71335.259 & No \\
        0.9 & -71160.774 & No \\
        \hline
    \end{tabular}
    \end{center}
    In particular, the initial guess of $0.7$ produces
    \[
        \hat{f}_{NM}=(0,\;0,\;0,\;0,\;1),
    \]
    which yields a substantially lower log-likelihood. Only one
    of the five initial guesses satisfies the convergence criteria.
    This suggests that Nelder-Mead can stagnate at inferior solutions,
    particularly when the simplex approaches the boundaries of the
    parameter space.

    \item[(c)]
    We now combine the two methods by first performing 30 iterations
    of Nelder-Mead and then using the resulting coefficients as the
    initial guess for Newton-Raphson.

    Starting from $f^{(0)}=(0.5,0.5,0.5,0.5,0.5)$, we obtain
    \[
        \hat{f}_{Hybrid}
        = (0,\;0,\;0.004654,\;0.071000,\;0.037850),
    \]
    with
    \[
        \mathcal{L}(Y\mid\hat{\theta}_{Hybrid})
        = -71160.713.
    \]
    The projected gradient error is $7.105\times10^{-13}$.

    Repeating the estimation over the five initial guesses considered
    above, we find that the hybrid method converges to the same
    coefficients and maximized log-likelihood in every case.
    Thus, the hybrid method is considerably less sensitive to the
    initial guesses than Nelder-Mead alone.

    \item[(d)]
    The results suggest that Newton-Raphson and the hybrid method
    perform better than Nelder-Mead for this estimation problem.
    Both methods converge to the same estimates across all five
    initial guesses, whereas Nelder-Mead converges for only one.

    Moreover, since the Bernstein hazard is linear in $f$, and the
    Bernoulli log-likelihood is concave in the replacement probability,
    the replacement log-likelihood is concave in $f$. Consequently,
    the differences in the estimates obtained with Nelder-Mead
    reflect numerical convergence problems rather than multiple
    local maxima.

    The estimated first two coefficients are zero, indicating that
    the solution lies on the boundary of the parameter space.
    This helps explain why the constrained Newton-Raphson method
    performs better than the simplex-based Nelder-Mead method.

    Finally, relative to Question 1, the maximized log-likelihood
    decreases from $-71146.863$ to $-71160.713$, a difference of
    approximately $13.850$. The Bernstein specification uses only
    five parameters to characterize the replacement hazard, rather
    than estimating 40 separate mileage-bin hazards. Thus, it
    provides a substantially more parsimonious representation
    of the replacement policy.
\end{enumerate}



\end{document}